% GENERATED FILE. Edit canonical lessons, then run npm run generate:books.
% canonical-source-hash: fnv1a64:2786694c80c8a2b4
% canonical-lessons: PA-R168-date-label-again, PA-R168-month-map-again, PA-R168-date-format-again

\chapter{Return to the Date Line}
\label{ch:pa-date-form-r4}
\hlchaptermodality{168}

\begin{chapteropening}
\textbf{By the end of this chapter:} \emph{I can retrieve three already-taught pieces of a fictional Punjabi date form after a long gap, one at a time, without claiming scored A1 writing or complete mock credit.}

Recall the date label, two taught month positions, and the printed day/month/year order without introducing new Punjabi.
\end{chapteropening}

\section[\texorpdfstring{\pa{ਤਾਰੀਖ਼} (tārīkh)}{ਤਾਰੀਖ਼ (tārīkh)}]{The old date label, once more}

\label{lesson:PA-R168-date-label-again}



\begin{quote}
Look at \textbf{\pa{ਤਾਰੀਖ਼}:} (\emph{tārīkh}). Recall that it labels the \textbf{date} line, not an age or phone line. The blank after it is for a date.
\end{quote}

\subsection*{Your turn}
Trace \textbf{\pa{ਤਾਰੀਖ਼}:} with a finger, then copy only the label and its colon once. Leave the value blank. Compare one mark at a time with the visible model.

\begin{tcolorbox}[breakable,colback=teal!4,colframe=teal!35!black,title={Before you move on}]
Cover the model. Which familiar label asks for a date? Write \textbf{\pa{ਤਾਰੀਖ਼}:}, then uncover and repair one mark if needed. This is retrieval practice, not a scored form submission.
\end{tcolorbox}
\section[\texorpdfstring{\pa{੦੧} · \pa{੦੨} (January …)}{੦੧ · ੦੨ (January …)}]{Two old month positions}

\label{lesson:PA-R168-month-map-again}



\begin{quote}
Look at \textbf{\pa{੦੧}} for January (\textbf{\pa{ਜਨਵਰੀ}}) and \textbf{\pa{੦੨}} for February (\textbf{\pa{ਫ਼ਰਵਰੀ}}). The first place is zero in both; the second changes.
\end{quote}

\subsection*{Your turn}
With both codes visible, copy \textbf{\pa{੦੧}} under January and \textbf{\pa{੦੨}} under February. Say the month, cover the code, and try its two digits once. Uncover and repair only the differing digit.

\begin{tcolorbox}[breakable,colback=teal!4,colframe=teal!35!black,title={Before you move on}]
From the old English cue \textbf{February}, write its month position \textbf{\pa{੦੨}}. Pause before looking. These are month positions, not new Punjabi words.
\end{tcolorbox}
\section[\_\_/\_\_/\_\_\_\_ (day/month/year)]{The old printed date order}

\label{lesson:PA-R168-date-format-again}



\begin{quote}
Look at the printed order from left to right: \textbf{day}, \textbf{month}, \textbf{year}. The familiar model \textbf{\_\_/\_\_/\_\_\_\_} has two places, two places, then four places. Slashes separate the boxes; they are not digits.
\end{quote}

\subsection*{Your turn}
Copy only the empty frame \textbf{\_\_/\_\_/\_\_\_\_}. Label its boxes in English: day, month, year. Keep the two slash marks visible; do not fill a real date.

\begin{tcolorbox}[breakable,colback=teal!4,colframe=teal!35!black,title={Before you move on}]
Cover the model. Draw the empty frame and point to the \textbf{middle month} box. Uncover and check its order and two separators. This is a no-stakes retrieval, not independent scored writing.
\end{tcolorbox}
